\section{Problem Formulation}

Let a flight history be a sequence of observations $X = \{x_1, x_2, \dots, x_T\}$, where each $x_t$ contains latitude, longitude, altitude, speed, heading, vertical rate, and derived temporal and mask features. In our experiments, the history spans 30 minutes and is resampled at 5-minute intervals, producing a short but regular input sequence with explicit observation and gap indicators.

For each forecast horizon $h \in \{10, 30, 60, 120\}$ minutes, the goal is to predict: (1) a probability distribution over target H3 cells $p_h(r \mid X)$, (2) a coarse parent distribution used for hierarchical decoding, and (3) an airborne/landed status distribution.

The target is not a single coordinate but a future regional distribution. This is motivated by the fact that long-horizon future movement is multimodal even under clean observations, and the uncertainty is amplified further in open operational data.

We use a hierarchical output space. For each horizon, the fine target H3 resolution is paired with a coarser parent resolution. The model first predicts the coarse parent region and then predicts the child cell conditional on that parent. This reduces output complexity and matches the hierarchical structure of H3.

Given predicted probabilities, we evaluate both classification quality and candidate-region quality. The main probabilistic metrics are NLL, Brier score, and ECE. For candidate regions, the headline metrics are cover at a fixed per-horizon cell budget and candidate union area, while candidate union cover and candidate union cell count are retained as secondary diagnostics. This evaluation design reflects the real deployment question more directly than top-k accuracy alone: a useful regional predictor should be accurate, calibrated, and spatially efficient, and broad coverage should not be rewarded unless it remains compact.
